% Pursuing Stacks (cotes 134-2 to 134-8): archive cote 134-6, site page 62 (PDF page 63), typescript page 380.
% typescript. Draft compared: ps5.tex:4506-4572.
% AI-transcribed by Claude Opus 5.5 (claude-opus-5-5, adaptive thinking, effort high), Message Batches API, batch ps3 (msgbatch_01XdFmuVVBAcgc7EFxdtCqsj), 2026-11-12; reviewed against both supplied scan halves; see CHECK_DECISIONS.tsv.
\scanpage{134-6/62}
\mspage{380}

\inserted{abelianization of} homotopy types modelized by a var\over[typed]{y}{i}able object \inserted[typed]{$X$} in $A^\wedge$, and b) we want that
the chain complex in (Ab) expressing abelianization of X, should be
expressible in terms of the "tautological abelianization" $\over{\mathrm{Wh}}{\mathrm{Wh}}_A(\over[typed]{\unsure{X}}{X}) = \underline{Z}^{(X)}$
of X , \inserted{itself} by a formula \inserted{moreover} wh\over{c}{i}\over{i}{c}h should make sense functorially with respect
to an arbitrary \over[typed]{e}{a}belian presheaf \unsure{i.e.} \inserted{an} object F in $A^\wedge_{ab}$ .

T\(\underset{\text{h}}{\text{h}}\)e main fact I have in view here is that whenever the chain comple\unsure{x}
L. in Addinf(B) is endowed with an augmentation
\[ (10) \qquad L_{\bullet} \hand{\longrightarrow} \underline{Z}_B \]
turning it into a resolution of $\underline{Z}_B$ i.e. into a (quasi special) \hand{\uline{\textcolor{black}{integra-}}}%
\hand{\uline{\textcolor{black}{tor}}} for B, then ipso facto L. is an abelianizator for A, the commuta-%
tion isomorphism \hand{$\lambda$} being canonicall\over[typed]{\unsure{a}}{y} defined by the augmentation (10).

Some comments, before proc\inserted[typed]{e}eding to a proof. Presumably, the
converse of our statement holds too - namely that the natural functor
we'll get from quasi-special in\over[typed]{e}{t}\over[typed]{t}{e}\over{\illegible{1}}{g}rators for B to abelianizators \inserted{$(L_{\bullet},\lambda)$} for A
is an equivalence \inserted[typed]{(even an isomorphism !)} between the rel\over[typed]{i}{e}vant categories. I do'nt feel like
pursuing this - \xout{more} the more relevant fact here, whether or not a con-%
verse as contemplated holds, is that we can pin down at any rate a
special class of \xout{inte} abelianizators for A, namely those which come from
\inserted{(quasi special)} integrators for B, and these abelianizators are defined up to \inserted[typed]{chain} homotopism
in Addinf(B). In this sense, \hand{\uline{\textcolor{black}{we get an existence and unicity statement}}}
\inserted{\uline{for abelianiz-ators in $A$},} as strong as we possibly could hope for. In practical terms, it would
seem, an \hand{\uline{\textcolor{black}{abelianizator for A}}} \hand{\uline{\textcolor{black}{will be no more no less than}}} \hand{\uline{\textcolor{black}{just a\xout{n}}}} \inserted[typed]{\hand{(}quasi-special\hand{)}} \hand{\uline{\textcolor{black}{inte-}}}%
\hand{\uline{\textcolor{black}{grator for B}}}, namely \hand{\uline{\textcolor{black}{a projective resolution of}}} $\underline{Z}_B$ in $B^\wedge_{ab}$ , \xout{having}
\hand{\uline{\textcolor{black}{whose components satisfy a mild extra assumption besides being projective}}}.

Here I am struck by a slight discrepancy in terminology, as we
would rather have a correspondance
\diagram{
\[
\hand{\left\{\textcolor{black}{
\begin{array}{l}
\text{integrators for B} \hand{\longrightarrow} \text{\over{\unsure{A}}{a}belianizators for A} \\[1ex]
\text{quasi-spe\xout{a}cial int\hand{\ensuremath{{}^{s}}}. for B} \hand{\longrightarrow} \text{quasi-special abelian\hand{\ensuremath{{}^{s}}}. for A ,}
\end{array}
}\right.}
\]
}
and the same for \xout{x}"special" integrators and abelianizators. As I still
feel that the general \xout{name of} appelation of an \hand{``}integrator\hand{''} for \hand{\uline{\textcolor{black}{any}}} projec-%
tive resolution of $\underline{Z}_B$ is adequate (without insisting that the components
shou\over{d}{l}\over{d}{d} be in Addinf(B)), this kind of forces upon us to extend according-%
ly \inserted{still} the notion of an abelianizator for A. This does make a sense, using
the pairing (6) (which we had dismissed as an "irrelevant nicety\xout{\unsure{"}} for
the time being" !) , and the corresponding equivalence

\margin{62}
